\section{Related Work}
\label{sec:related-work}

\paragraph{Coverage-guided fuzzing.}
American Fuzzy Lop~\cite{zalewski2014afl} and its successor AFL++~\cite{fioraldi2020aflplusplus}
popularized mutation-based, coverage-guided fuzzing: a target is
instrumented to report which code ran, and the fuzzer retains and mutates
inputs that reach new coverage. Man\`es et al.'s survey~\cite{manes2021fuzzing}
gives a taxonomy of this and related approaches. Coverage-guided fuzzers are
highly effective when instrumentation is available, but instrumentation
itself is exactly what a memory-safe, JIT/AOT-compiled runtime such as
Dart's makes both less necessary (there is no memory-safety class of bug for
a sanitizer to catch) and less available in the form these tools assume;
none of this family targets the failure modes this paper is concerned with.

\paragraph{Grammar-based and generation-based fuzzing.}
Generating inputs directly from a grammar, rather than mutating bytes, is a
much older idea: Purdom's sentence generator~\cite{purdom1972sentence}
produces short strings that exercise every grammar production, originally
for testing and debugging parsers themselves. More recent tools combine a
grammar with coverage feedback: Superion~\cite{wang2019superion} extends a
greybox fuzzer with grammar-aware mutation operators so mutated inputs
remain syntactically valid, and Nautilus~\cite{aschermann2019nautilus}
similarly pairs a grammar-based generator with coverage-guided mutation to
reach deeper parser states. Both still assume coverage instrumentation is
available; the pipeline this paper extends removes that assumption
entirely.

\paragraph{Property-based testing.}
QuickCheck~\cite{claessen2000quickcheck} introduced property-based testing:
a developer writes generators for values of a given shape and a property
that should hold for all of them, and the tool searches for a
counterexample. Hypothesis~\cite{maciver2019hypothesis}, which the pipeline
this paper reuses builds directly on, is a Python implementation of the
same idea, with composable strategies expressive enough to encode both a
formal grammar (our RQ1 generator) and an adversarial perturbation catalog
over a target-specific schema (our RQ2 generator, Section~\ref{sec:methodology}).

\paragraph{LLMs for fuzzing and test generation.}
LLMs have recently been used to generate or mutate fuzzing inputs directly.
TitanFuzz~\cite{deng2023titanfuzz} uses an LLM to generate and mutate
programs for fuzzing deep-learning library APIs, relying on the LLM's prior
exposure to real usage examples of those APIs rather than a formal grammar.
Fuzz4All~\cite{xia2024fuzz4all} generalizes this to arbitrary
languages/systems by using an LLM both to synthesize a distilled input
prompt from documentation and examples and to generate and mutate inputs
from it. CodaMOSA~\cite{lemieux2023codamosa} invokes an LLM only after
search-based test generation's coverage plateaus, asking it for example
test cases to escape the plateau -- i.e., the LLM is invoked \emph{because}
coverage feedback is available and has stalled. This paper's setting differs
along the same axis as its predecessor~\cite{mansy2026agentic}: no coverage
signal is available at any point, and the LLM's job is to revise a reusable,
executable \emph{generator} against feedback a black-box harness can
actually produce, not to emit one-off inputs.

\paragraph{Differential testing.}
Running multiple independent implementations of the same specification
against identical inputs and comparing their outputs, rather than checking
each against a hand-written oracle, dates to McKeeman's differential
testing~\cite{mckeeman1998differential}. This idea is what makes RQ2's
oracle possible at all in a setting with no crash to detect: four
independently-implemented Dart JSON deserialization paths sharing one
formal schema give a comparison target for free, the same way McKeeman
compared independent compiler implementations against each other rather
than against a formal semantics. Where classic differential testing
compares whole programs given the same input, our four paths share almost
their entire input pipeline (Section~\ref{sec:methodology}'s shared
JSON-syntax gate) and differ only in schema-level decoding, which is a
narrower and, we find, more failure-prone surface than differential testing
is usually applied to.

\paragraph{Positioning.} This paper is a direct empirical follow-up to
\cite{mansy2026agentic}, which evaluated the identical refinement-loop
mechanism on two independently implemented C JSON parsers under a
sanitizer-based crash oracle and named a memory-safe target as future work
along a materially different generalization axis. We take up exactly that
axis: same loop, same sandboxing discipline, same reproducibility tooling
where none of it is C- or format-specific, retargeted at a language where
the entire premise of a crash oracle does not hold, using differential
divergence~\cite{mckeeman1998differential} and ground-truth numeric
round-tripping as replacements. Unlike TitanFuzz, Fuzz4All, and CodaMOSA,
the LLM here never sees coverage data, execution traces, or usage examples --
only a schema description and the same three-signal proxy discipline (here,
a divergence score in place of acceptance rate) a human maintaining a
Hypothesis strategy by hand would have to reason about manually.
